\documentclass[11pt,a4paper]{article}
\usepackage{fontspec}
% Liberation Sans tem as medidas da Arial; sem ela, usa a TeX Gyre Heros
\IfFontExistsTF{Liberation Sans}{\setmainfont{Liberation Sans}}{\setmainfont{TeX Gyre Heros}}
\IfFontExistsTF{Liberation Mono}{\setmonofont{Liberation Mono}}{\setmonofont{TeX Gyre Cursor}}
\usepackage{polyglossia}
\setdefaultlanguage[variant=brazilian]{portuguese}
\usepackage[left=3cm,right=2cm,top=3cm,bottom=2cm,headsep=0.8cm]{geometry}
\usepackage{setspace}
\setstretch{1.35}
\setlength{\parindent}{0pt}
\setlength{\emergencystretch}{2.5em}
\setlength{\parskip}{6pt plus 1pt}
\usepackage{graphicx}
\usepackage[table]{xcolor}
\usepackage{array,booktabs,tabularx,multirow,makecell}
\usepackage{fancyhdr}
\usepackage{titlesec}
\usepackage{tocloft}
\usepackage{caption}
\usepackage{float}
\usepackage[section,above]{placeins}
\usepackage{flafter}
\usepackage{fancyvrb}
\usepackage{multicol}
\setlength{\columnsep}{0.8cm}
\usepackage{amsmath}
\usepackage{tikz}
\usetikzlibrary{positioning,calc,arrows.meta,fit,backgrounds,shapes.geometric,decorations.pathreplacing}
\usepackage[hidelinks]{hyperref}
\usepackage{xurl}
\urlstyle{same}
\hypersetup{pdftitle={Lab de configuração básica de rede com VLANs e VLSM},
  pdfauthor={Gabriel André de Lima Silva},
  pdfsubject={Atividade de Redes I, Professor Fred Sauer}}

\definecolor{hostblue}{HTML}{2F5DA8}
\definecolor{vprof}{HTML}{1F5FAD}
\definecolor{valun}{HTML}{2E7D32}
\definecolor{vvisi}{HTML}{B85C00}
\definecolor{gridgray}{gray}{0.72}
\newcommand{\hb}[1]{\textcolor{hostblue}{#1}}

% cabecalho: numero da pagina no canto superior direito
\pagestyle{fancy}
\fancyhf{}
\fancyhead[R]{\small\thepage}
\renewcommand{\headrulewidth}{0pt}

% titulos
\titleformat{\section}{\normalfont\bfseries}{}{0pt}{}
\titleformat{\subsection}{\normalfont\bfseries}{}{0pt}{}
\titleformat{\subsubsection}{\normalfont\bfseries}{}{0pt}{}
\titlespacing*{\section}{0pt}{14pt}{6pt}
\titlespacing*{\subsection}{0pt}{10pt}{4pt}
\titlespacing*{\subsubsection}{0pt}{8pt}{2pt}
\newcommand{\secao}[1]{\section*{#1}\phantomsection\addcontentsline{toc}{section}{#1}}
\newcommand{\subsecao}[1]{\subsection*{#1}\phantomsection\addcontentsline{toc}{subsection}{#1}}

% sumario
\renewcommand{\cfttoctitlefont}{\hfill\normalsize\bfseries}
\renewcommand{\cftaftertoctitle}{\hfill}
\setlength{\cftbeforetoctitleskip}{0pt}
\setlength{\cftaftertoctitleskip}{14pt}
\renewcommand{\cftsecfont}{\small\bfseries}
\renewcommand{\cftsecpagefont}{\small\bfseries}
\renewcommand{\cftsubsecfont}{\small}
\renewcommand{\cftsubsecpagefont}{\small}
\renewcommand{\cftsecleader}{\bfseries\cftdotfill{\cftdotsep}}
\setlength{\cftbeforesecskip}{2pt}
\setlength{\cftsubsecindent}{1.2em}
\setlength{\cftsecindent}{0pt}

% legendas: "Tabela 1: ..." acima, fonte embaixo
\captionsetup{font=small,labelsep=colon,justification=raggedright,singlelinecheck=false,skip=3pt}
\captionsetup[table]{position=top}
\newcommand{\fonte}{\par\vspace{1pt}{\small Fonte: elaborado pelo autor.}\par}
\newcommand{\fontebase}[1]{\par\vspace{1pt}{\small Fonte: elaborado pelo autor com base em #1.}\par}
\newcommand{\fonteprint}{\par\vspace{1pt}{\small Fonte: captura de tela do Cisco Packet Tracer feita pelo autor.}\par}
\newcommand{\fonteadapt}{\par\vspace{1pt}{\small Fonte: adaptado de Sauer ([\emph{s. d.}]a).}\par}

% tabela de contas em binario (grade cinza, cabecalho sombreado)
\newenvironment{bintab}{%
  \par\vspace{2pt}\noindent\begingroup\small\arrayrulecolor{gridgray}\renewcommand{\arraystretch}{1.25}%
  \begin{tabular}{|l|l|l|}\hline
  \rowcolor{gray!14} & \multicolumn{1}{c|}{\bfseries Binário (rede | host)} & \multicolumn{1}{c|}{\bfseries Decimal}\\ \hline}
 {\end{tabular}\endgroup\par\vspace{6pt}}
\newenvironment{andtab}{%
  \par\vspace{2pt}\noindent\begingroup\small\renewcommand{\arraystretch}{1.15}\begin{tabular}{@{}l@{\hspace{1em}}l@{\hspace{1em}}l@{}}}
 {\end{tabular}\endgroup\par\vspace{6pt}}
\newenvironment{somatab}{%
  \par\vspace{2pt}\noindent\begingroup\small\renewcommand{\arraystretch}{1.15}\begin{tabular}{@{}l@{\hspace{0.6em}}r@{\hspace{1em}}l@{}}}
 {\end{tabular}\endgroup\par\vspace{6pt}}

% blocos de comandos
\DefineVerbatimEnvironment{cmd}{Verbatim}{fontsize=\small,xleftmargin=1.2em,frame=leftline,framerule=0.6pt,rulecolor=\color{gray!60},framesep=8pt,samepage=true,baselinestretch=1.08}
\DefineVerbatimEnvironment{conf}{Verbatim}{fontsize=\scriptsize,baselinestretch=1.0}
\DefineVerbatimEnvironment{saida}{Verbatim}{fontsize=\footnotesize,xleftmargin=0.4em,frame=single,rulecolor=\color{gray!50},framesep=6pt,samepage=true,baselinestretch=1.0}

\newcolumntype{C}[1]{>{\centering\arraybackslash}p{#1}}
\setlength{\cftfigindent}{0pt}

% icones e estilos comuns aos diagramas
\tikzset{
  sw/.pic={\draw[fill=white, rounded corners=2pt] (-0.55,-0.22) rectangle (0.55,0.22);
           \draw[-{Stealth[length=3pt]}] (-0.32,0.07) -- (0.32,0.07);
           \draw[-{Stealth[length=3pt]}] (0.32,-0.07) -- (-0.32,-0.07);},
  rt/.pic={\draw[fill=white] (0,0) circle (0.36);
           \draw[{Stealth[length=3pt]}-{Stealth[length=3pt]}] (-0.22,-0.22) -- (0.22,0.22);
           \draw[{Stealth[length=3pt]}-{Stealth[length=3pt]}] (-0.22,0.22) -- (0.22,-0.22);},
  pc/.pic={\draw[fill=white, rounded corners=1pt] (-0.32,-0.12) rectangle (0.32,0.30);
           \draw (0,-0.12) -- (0,-0.22); \draw (-0.16,-0.24) -- (0.16,-0.24);},
  porta/.style={font=\scriptsize, text=black!80, fill=white, inner sep=0.6pt},
  info/.style={draw, rounded corners=1.5pt, align=left, font=\scriptsize, inner sep=3pt,
               line width=0.6pt, text width=2.65cm},
}
% tronco 802.1Q: tres linhas paralelas, uma por VLAN
\newcommand{\tronco}[2]{%
  \draw[vvisi, line width=0.9pt] ($(#1)!1.8pt!90:(#2)$) -- ($(#2)!1.8pt!-90:(#1)$);
  \draw[valun, line width=0.9pt] (#1) -- (#2);
  \draw[vprof, line width=0.9pt] ($(#1)!1.8pt!-90:(#2)$) -- ($(#2)!1.8pt!90:(#1)$);}

% listas de figuras e de tabelas
\tocloftpagestyle{empty}
\renewcommand{\cftloftitlefont}{\hfill\normalsize\bfseries}
\renewcommand{\cftafterloftitle}{\hfill}
\renewcommand{\cftlottitlefont}{\hfill\normalsize\bfseries}
\renewcommand{\cftafterlottitle}{\hfill}
\setlength{\cftbeforeloftitleskip}{0pt}
\setlength{\cftafterloftitleskip}{14pt}
\setlength{\cftbeforelottitleskip}{0pt}
\setlength{\cftafterlottitleskip}{14pt}
\renewcommand{\cftfigpresnum}{Figura~}
\renewcommand{\cftfigaftersnum}{:}
\setlength{\cftfignumwidth}{5.6em}
\renewcommand{\cfttabpresnum}{Tabela~}
\renewcommand{\cfttabaftersnum}{:}
\setlength{\cfttabnumwidth}{5.6em}
\setlength{\cftbeforefigskip}{3pt}
\setlength{\cftbeforetabskip}{3pt}
